\section{Modeling choices explained}
\label{app:modeling}

This appendix explains why the model of Section~\ref{sec:model} has the shape it
does. It is written as a deliberate explanatory layer: it does not prove
anything, and each choice is labeled by its origin --- historical/source
mandated, mathematical idealization, operational approximation, or
project-level extension. Where an external treatment explains a convention
better than we can inline, we cite it rather than reproduce background.

\subsection{Why the genome is circular}

\paragraph{Source origin.}
Shomorony et al.\ work with a circular sequence ``for ease of exposition \dots
[to] avoid edge effects,'' so that every start position has a length-$L$ window
\cite[\S2]{shomorony2016}. Medvedev and Brudno's exact likelihood is likewise
stated for a ``circular genome'' \cite[\S6.1]{medvedev2009}.

\paragraph{What it buys.}
Circularity removes the two ends: the window map $t\mapsto W_t(D)$ is defined at
every position, the spectrum has total $|D|$ with no boundary correction, and the
window-support graph of Definition~\ref{def:graph} is balanced and strongly
connected because it is generated by a single closed walk. Several proofs would
need extra boundary terms without this.

\paragraph{What it approximates.}
A real chromosome is linear; the circular model is appropriate when the genome is
genuinely circular (many bacterial genomes) or when reads near the ends are rare
relative to the genome length, so edge effects are negligible. It is an
idealization, not a claim that genomes are circles.

\paragraph{Mandatory consequence.}
Cyclic-shift equivalence is forced by circularity (Remark~\ref{rem:cyclic}).
A formalization that models a circle but compares literal strings is internally
inconsistent.

\subsection{Why reads are fixed-length, error-free, and uniformly sampled}

\paragraph{Source origin.}
The 2016 exposition draws each of $N$ reads independently and uniformly from the
$G$ circular length-$L$ substrings of $s$ \cite[\S2]{shomorony2016}. Bresler et
al.\ use the same noiseless uniform model with bridging \cite{bresler2013}.

\paragraph{Why errors can be abstracted away in the source theory.}
Shomorony et al.\ deliberately separate an error-free information-theoretic phase
from sequencing-error handling; their algorithm uses exact overlaps, and their
double-strand implementation adds reverse-complement reads as preprocessing
\cite[\S4.1]{shomorony2016}. In an information-feasibility analysis this is a
modeling choice: it isolates the repeat-limited obstructions that determine the
information boundary. What the abstraction omits is real: sequencing errors
create spurious overlaps, turn unique repeats ambiguous, and interact with
coverage in ways a fixed $L$ cannot express. We therefore do not claim the
error-free model describes production assemblers.

\paragraph{Why uniform and independent.}
Uniformity over start positions makes the exact likelihood a multinomial with
probabilities $d_i/N(D)$; a non-uniform or position-dependent coverage profile
would destroy that conjugacy and change the objective. Independence makes the
joint likelihood a product, which is what enables both the exact multinomial and
the separable approximation. Neither is a biological law; both are the minimal
statistical assumptions under which the source's likelihood has the stated form.

\subsection{Why multiplicity is retained}

The likelihood is a function of the observed read-type \emph{counts}
$x(w)$, not of the set of distinct read strings. This is source-mandated: the
exact objective is the multinomial
$\frac{n!}{\prod_w x(w)!}\prod_w(d_w/N(D))^{x(w)}$. Collapsing reads to a set
would discard the frequencies that the objective scores and would make the
finite-sampling failure of Example~\ref{ex:frequencies} invisible. Keeping the
latent true placements separate from the observable counts is equally important:
bridging is a property of $(S,R)$, where $R$ records placements, while the
likelihood estimator may not see them.

\subsection{Oriented reads versus reverse-complement molecules}

\paragraph{Two genuinely different sources.}
Shomorony's theory uses oriented length-$L$ substrings; reverse complement enters
only as experimental preprocessing that \emph{adds} reads \cite[\S4.1]{shomorony2016}.
Medvedev and Brudno instead model a read as a DNA \emph{molecule}, an unordered
pair of reverse-complementary strands, represented once in the graph
\cite[\S3.1,\S4.1]{medvedev2009}. This is not notation: the two models have
different read-type spaces.

\paragraph{Consequences.}
Under oriented types, reverse-complementing a candidate relabels its type counts
by $t\mapsto\rc(t)$, and the candidate and its reverse complement need not tie.
In fact they tie only when the observed count vector is reverse-complement
symmetric, which fails in general, e.g.\ $D=AAAT$ at $L=1$ has $A$-count $3$
while $\rc(D)=TTTA$ has $A$-count $1$. Under molecule classes, a candidate and
its reverse complement \emph{always} tie because class counts are
reverse-complement invariant. Hence the read-type convention and the genome
equivalence are coupled (Section~\ref{sec:open}): molecule types force dihedral
equivalence, oriented types force only cyclic shift. The internal tension in
Medvedev and Brudno's Section~6.1 --- the text writes ``$4^k$ such variables''
while its vocabulary is $k$-molecules --- must be resolved per formalization
rather than glossed.

\subsection{Candidate length, known genome size, and same-length restriction}

These are three distinct assumptions, and conflating them is the most common
source-faithfulness error:

\begin{enumerate}[label=(\alph*),leftmargin=*]
  \item the likelihood may have a \emph{candidate-intrinsic} length $N(D)$
        (exact multinomial);
  \item it may be supplied an \emph{external} genome-size parameter $N$, used in
        the per-type binomial marginals (the Section~6.1 approximation);
  \item the \emph{candidate class} may be restricted to $|D|=G$.
\end{enumerate}
(a) and (b) are different objectives; (b) does not imply (c). Medvedev and
Brudno say ``we assume that the genome size is known''
\cite[\S6.1]{medvedev2009}, and a contemporaneous formulation fixes the assembly
length to the known genome length and justifies it by scale-degeneracy
\cite{ghodsi2016}. But the 2016 sentence itself fixes neither, so a same-length
theorem is a restricted result unless a source argument fixes the restriction.
The fixed-length witnesses of Section~\ref{sec:finite} are same-length precisely
to be robust: by Lemma~\ref{lem:negative-transfer} they refute the statement over
all circular candidates as well.

\subsection{Finite sample versus population}

A finite sample is what a sequencing experiment provides; a population is the
distribution it is drawn from. The two answer different questions. In the finite
regime the empirical frequencies fluctuate, and a wrong candidate can win
(Example~\ref{ex:frequencies}); the question is whether bridging is strong enough
to prevent this anyway, which Section~\ref{sec:finite} answers in the negative
for precise readings. In the population regime the fluctuation is gone, so the
maximizer statement is automatic and only identifiability remains
(Section~\ref{sec:population}). The population model is a project-level
idealization used to isolate the identifiability question; it is not the
historical finite-data model, and it does not settle it.

\subsection{What this paper deliberately does not model}

To keep the main narrative readable we excluded several directions that the
repository has not studied in depth: linear genomes with genuine end effects;
variable-length reads; noisy or adversarial reads beyond the source's
abstraction; and the Section~6.2 flow's non-contiguous assemblies as
sequence-valued candidates. We list them here rather than silently omitting
them, and we cite the adjacent literature on variable-length reads
\cite{hui2016} and safe/complete assembly \cite{tomescu2016,rahman2022} for
readers who want those directions.
